% !TeX root = ../main.tex
% ---------------------------------------------------------------------------
%  Resilience of Underground Infrastructure
%  Added from the author's study report on sustainability and resilience of
%  underground infrastructure, which supplies the functionality-over-time
%  framing, the resilience assessment sequence and the restoration decision
%  logic used here.
%  Chapter placed in Part VI so that the hazard chapters that precede it feed
%  a common lifecycle metric rather than remaining separate design checks.
% ---------------------------------------------------------------------------
\chapter{Resilience of Underground Infrastructure}
\label{ch:resilience}

The hazard chapters that precede this one each treat a single threat: seismic
shaking, water inrush, squeezing, deterioration. Designed against
independently, they produce a set of unrelated capacity checks. What an asset
owner needs instead is a common measure: how much service function is lost when
something happens, and how quickly it returns.

That measure is resilience, and it is expressed through functionality over
time rather than through hazard intensity. An underground asset begins its life
delivering a required level of service. Ageing, creep, deterioration and
leakage erode that level gradually; earthquakes, blasting, pressure transients
and construction events reduce it abruptly. Inspection, repair, strengthening,
grouting and relining restore it. Resilience is the shape of that curve.

\section{Functionality as the design measure}
\label{sec:functionality-measure}

\subsection{Why hazard intensity is not enough}
\label{sec:beyond-hazard}
\scopenote{A design that survives the design event may still be unusable for
months. Capacity checks answer whether the structure stands; they do not
answer whether the service continues, which is what the owner and the public
experience.}

\subsection{Defining the service function}
\label{sec:service-function}
\scopenote{Expressing function in terms the asset actually delivers ---
traffic capacity, flow rate, safe occupancy, generating availability --- and
setting the threshold below which the asset is considered out of service.}

\section{Deterioration and hazard scenarios}
\label{sec:deterioration-and-hazards}

\subsection{Baseline life-cycle deterioration}
\label{sec:baseline-deterioration}
\scopenote{The slow decline against which events are superimposed: ageing,
creep, corrosion, leaching, leakage and progressive lining damage. Mechanisms
in Chapter~\ref{ch:asset-management}.}

\subsection{Discrete hazard events}
\label{sec:discrete-hazards}
\scopenote{Earthquake (Chapter~\ref{ch:seismic-design}), blasting and
construction disturbance, operational pressure transients, water inrush
(Chapter~\ref{ch:construction-hazards}), impact and fire. Selecting scenarios
proportionate to consequence.}

\subsection{Multi-hazard interaction}
\label{sec:multi-hazard}
\scopenote{Deterioration raises vulnerability to events, and events accelerate
deterioration. Treating the two independently understates lifetime risk.}

\section{Resilience assessment}
\label{sec:resilience-assessment}

\subsection{The assessment sequence}
\label{sec:resilience-sequence}
A workable assessment proceeds in eight steps:

\begin{enumerate}[nosep]
  \item define the service function as a functionality curve over time;
  \item establish baseline performance under ageing and deterioration;
  \item apply the hazard scenarios;
  \item quantify functionality loss immediately after each event, giving
        vulnerability and robustness;
  \item estimate residual functionality;
  \item estimate delay time before restoration begins, and restoration
        duration;
  \item evaluate the recovery trajectory;
  \item compute a resilience index and compare intervention options.
\end{enumerate}

\subsection{Resilience metrics}
\label{sec:resilience-metrics}
\scopenote{Initial functionality, hazard occurrence time, immediate loss,
residual functionality, delay, recovery trajectory, final restored
functionality, and the area under the functionality--time curve. Recovery
shape matters as much as recovery duration.}

\subsection{Uncertainty in the resilience estimate}
\label{sec:resilience-uncertainty}
\scopenote{Uncertainty enters through hazard intensity, damage state,
detection delay and recovery rate. Reporting a single index without those
bands overstates what the assessment establishes. Methods in
Chapter~\ref{ch:reliability-codes}.}

\section{Restoration as a design decision}
\label{sec:restoration}

\subsection{Comparing intervention options}
\label{sec:intervention-comparison}
\scopenote{Restoration schemes compared on functionality recovered, time,
cost, constructability, risk reduction and carbon impact --- not on repair
cost alone. Repair options are design decisions with service and carbon
consequences.}

\subsection{Designing for recoverability}
\label{sec:design-for-recovery}
\scopenote{Access for inspection and repair, replaceable elements, redundancy
and sacrificial detailing at fault crossings. Decisions taken at design stage
that determine how fast the asset can be brought back.}

\section{Closing the loop}
\label{sec:resilience-loop}

\subsection{Inspection and post-intervention feedback}
\label{sec:resilience-feedback}
\scopenote{Condition data and post-repair performance feed back into the
digital twin of Chapter~\ref{ch:digital-underground}, updating the baseline
for the next assessment cycle.}

\subsection{From reactive repair to forecast-based intervention}
\label{sec:predictive-intervention}
\scopenote{Surface observation, internal inspection and multi-physics
modelling together support intervention before functionality is lost, which is
the practical objective of predictive maintenance
(Chapter~\ref{ch:asset-management}).}
